% ---------------------------------------------------------------------------
% Appendix: static sink-probe validity.
% Patterns are copied VERBATIM from Benchmarks/extraction/extract_tools.py
% (the SINKS dict). If that file changes, this appendix is stale.
% Rates are from Benchmarks/extraction/results.jsonl, n=597 Python / 807 Node.
% ---------------------------------------------------------------------------

\section{Validity of the Static Sink Probes}
\label{app:probes}

Table~\ref{tab:sink-families} reports, for each of ToFu's 21 introspection
tracers, the fraction of corpus repositories containing a static reference to
an API that the tracer instruments. Those fractions are produced by regular
expressions over repository source, and the expressions vary widely in how
specific they are. We document that variation here rather than suppress it,
because it determines what the table can and cannot support.

\paragraph{How the probes were built.}
ToFu's 21 introspection tracers are anchored in CWE: each oracle names the
weakness class it monitors (\eg{} CWE-502 for unsafe deserialization, CWE-918
for SSRF), and CWE serves as the reusable security knowledge behind the
sanitizer corpus (\S\ref{sec:design}). The \emph{static probes} used to
measure corpus coverage are a separate, later artefact and are not
taxonomy-derived: they are the inverse image of the instrumentation, one
regular expression per tracer, written to approximate ``would this tracer ever
have anything to observe in this repository.'' Each repository is cloned at depth
one; every file with a Python or JavaScript/TypeScript extension is scanned,
excluding \texttt{node\_modules}, virtual environments, build output and test
directories, and files over 900\,KB; a family is recorded as present if any
pattern matches anywhere in the repository. Matching is therefore at
repository granularity and is not confined to the bodies of registered tools.

\paragraph{Two regimes.}
The probes fall into two groups. The first targets a distinctive API whose
presence genuinely indicates the corresponding behaviour: \texttt{pickle} and
\texttt{yaml.load} for \textsc{Deserialization}, \texttt{verify=False} and
\texttt{rejectUnauthorized: false} for \textsc{SSLCertValidation},
\texttt{lxml} and \texttt{xml2js} for \textsc{XXE}. These probes report low
rates (2--5\%), and we regard their rates as informative.

The second group, listed in Table~\ref{tab:probe-validity}, matches ordinary
language constructs. \textsc{LogInjection} matches \texttt{print(}, which
appears in nearly every Python file; \textsc{FileAccess} matches any
\texttt{open(}; \textsc{ReDoS} matches any use of the \texttt{re} module.
A repository scoring positive for these families has not been shown to be at
risk---it has been shown to be written in Python or JavaScript. Their high
rates (16--85\%) are a property of the probe, not of the ecosystem.

\paragraph{Why the distinction matters.}
We use the coverage numbers for one purpose only: to argue that a tracer has
enough reachable surface in the corpus to be worth running. For that purpose
a low-specificity probe is still a valid \emph{upper bound}---a repository
with no \texttt{open(} anywhere certainly has no file-access defect---and the
nine families in Table~\ref{tab:probe-validity} clear it comfortably. We do
not use these nine families to stratify samples, to rank threat prevalence,
or to form denominators for detection rates, because a denominator built from
\texttt{print(} is not a population of at-risk servers.

\paragraph{Overlap.}
The probes are not disjoint. \textsc{BlockingEventLoopCall} shares
\texttt{readFileSync} with \textsc{FileAccess} and \texttt{execSync} with
\textsc{CommandExecution}, so a single Node file using both constructs
contributes to three families. Marginal rates in
Table~\ref{tab:sink-families} therefore cannot be summed, and no
joint-distribution claim should be read from them.

\paragraph{Unvalidated precision.}
We did not hand-label a sample, so the probes have no measured precision or
recall, and the coverage rates carry no error bar. We state the rates as
properties of the scan rather than estimates of a latent quantity.
Establishing precision would require a labelled sample per family and is
future work. Two further limitations follow from the method: static matching
cannot see dynamically constructed calls, and repository-granularity matching
attributes a sink to a repository even when it occurs in code no registered
tool can reach.

\begin{table}[t]
\centering
\caption{The nine low-specificity probes. Patterns are verbatim from the
  scanner; rates are the fraction of repositories matching
  ($n=597$ Python, $n=807$ Node/TS).}
\label{tab:probe-validity}
\footnotesize
\setlength{\tabcolsep}{4pt}
\begin{tabular}{@{}lrrp{2.6cm}p{2.9cm}@{}}
\toprule
\textbf{Family} & \textbf{Py} & \textbf{Node} & \textbf{Pattern} & \textbf{What it matches} \\
\midrule
LogInjection & 85.4 & 66.7 &
  \texttt{logging.} \texttt{logger.} \texttt{console.log} \texttt{print\textbackslash s*(} &
  Any file that prints or logs at all. \\
\addlinespace
SensitiveData\-Logging & 72.5 & 79.1 &
  \texttt{os.environ} \texttt{process.env} \texttt{getenv} \texttt{dotenv} &
  Any configuration read. No logging is implied by the match. \\
\addlinespace
FileAccess & 59.5 & 50.3 &
  \texttt{\textbackslash bopen\textbackslash s*(} \texttt{Path(} \texttt{pathlib} \texttt{fs.readFile} \texttt{readFileSync} &
  Any file open, including reading the program's own config. \\
\addlinespace
FileWrite & 34.8 & 47.6 &
  \texttt{fs.writeFile} \texttt{writeFileSync} \texttt{shutil.(copy|move)} \texttt{.write\textbackslash s*(} &
  Any \texttt{.write(} call, including \texttt{sys.stdout.write}. \\
\addlinespace
ReDoS & 31.0 & 20.0 &
  \texttt{re.compile} \texttt{re.(match|search|findall)} \texttt{new RegExp} &
  Any regular expression, regardless of whether it can backtrack. \\
\addlinespace
BlockingEvent\-LoopCall & 30.2 & 39.7 &
  \texttt{time.sleep} \texttt{execSync} \texttt{readFileSync} \texttt{sleep\textbackslash s*(} &
  Any sleep or sync I/O. Overlaps \textsc{FileAccess} and \textsc{CommandExecution}. \\
\addlinespace
LoopCondition & 27.8 & 16.9 &
  \texttt{while\textbackslash s+True} \texttt{while(\ true)} \texttt{for(;;)} &
  Any unbounded loop, including a normal server accept loop. \\
\addlinespace
Unrestricted\-IPBinding & 26.0 & 37.1 &
  \texttt{0.0.0.0} \texttt{host\textbackslash s*=\textbackslash s*"0} \texttt{listen\textbackslash s*(} &
  Any \texttt{listen(} call. Most matches are ordinary localhost servers. \\
\addlinespace
WeakPRNG & 16.4 & 24.0 &
  \texttt{\textbackslash brandom.} \texttt{Math.random} \texttt{randint} \texttt{choice(} &
  Any randomness, including \texttt{random.shuffle} on a list. \\
\bottomrule
\end{tabular}
\end{table}
